\section{Results Framework}
\label{sec:results-framework}

This section contains no empirical outcome. The word \texttt{pending} can be
replaced only by a generated artifact from an identified frozen run. Passing
\cref{eq:scientific-gate} is not required to report what a run observed. It is
required to label that run release-eligible and to authorize release claims.

\providecommand{\AFDReleaseGate}{pending}
\providecommand{\AFDRunStatus}{pending}
\providecommand{\AFDPredecessorRun}{pending}
\providecommand{\AFDRerunReason}{pending}
\providecommand{\AFDFalseAcceptanceNumerator}{pending}
\providecommand{\AFDFalseAcceptanceDenominator}{pending}
\providecommand{\AFDFalseAcceptanceRate}{pending}
\providecommand{\AFDFalseBlockingNumerator}{pending}
\providecommand{\AFDFalseBlockingDenominator}{pending}
\providecommand{\AFDFalseBlockingRate}{pending}
\providecommand{\AFDExpectedActionMismatches}{pending}
\providecommand{\AFDPipelineErrors}{pending}
\providecommand{\AFDReasonMismatches}{pending}
\providecommand{\AFDCanonicalGateFailures}{pending}
\providecommand{\AFDDifferentialFailures}{pending}
\providecommand{\AFDFidelityFailures}{pending}
\providecommand{\AFDFuzzFailures}{pending}
\providecommand{\AFDFaultFailures}{pending}
\providecommand{\AFDExactStableOutputs}{pending}
\providecommand{\AFDSemanticClosureOutputs}{pending}
\providecommand{\AFDPerformanceCI}{pending}
\providecommand{\AFDServiceRate}{pending}

\subsection{Run identity and release eligibility}
\label{subsec:results-identity}

Results begin with identity rather than eligibility. Run \AFDRunID uses commit
\AFDCommit, contains \AFDCorpusSamples corpus records, and has status
\AFDRunStatus. Its release-gate value is \AFDReleaseGate. A rerun cites
\AFDPredecessorRun and records \AFDRerunReason. The run remains in the registry
when the gate is false.

\begin{table}[tb]
\centering
\caption{Controlled-run identity and release eligibility. All entries are
generated placeholders rather than empirical claims.}
\label{tab:results-run}
\scriptsize
\begin{tabularx}{\textwidth}{L{0.38\textwidth}Y}
\toprule
Field & Generated value \\
\midrule
Run identifier & \AFDRunID \\
Source commit & \AFDCommit \\
Run status & \AFDRunStatus \\
Predecessor and rerun reason & \AFDPredecessorRun, \AFDRerunReason \\
Manifest records & \AFDCorpusSamples \\
Accepted outputs & \AFDAcceptedSamples \\
Blocked records & \AFDBlockedSamples \\
Pipeline errors & \AFDPipelineErrors \\
Scientific release gate & \AFDReleaseGate \\
\bottomrule
\end{tabularx}
\end{table}

\subsection{Action outcomes and finite-corpus rates}
\label{subsec:results-actions}

False acceptance and false blocking use separate denominators. The primary
false-blocking denominator contains only expected-clean S0 and S1 records. The
false-acceptance denominator contains every record whose pinned policy oracle
requires withholding. Expected-action mismatch is reported for the remaining
strata. A null machine value and the printed phrase ``not estimable'' represent
an empty denominator. Zero is reserved for an observed zero count over a
nonempty eligible set.

\begin{table}[tb]
\centering
\caption{Primary finite-corpus action outcomes. No population confidence
interval is attached to the purposive clustered corpus.}
\label{tab:results-binary}
\small
\begin{tabularx}{\textwidth}{L{0.34\textwidth}L{0.28\textwidth}Y}
\toprule
Endpoint & Events / eligible & Finite-corpus proportion \\
\midrule
False acceptance & \AFDFalseAcceptanceNumerator{} / \AFDFalseAcceptanceDenominator & \AFDFalseAcceptanceRate \\
Primary benign false blocking & \AFDFalseBlockingNumerator{} / \AFDFalseBlockingDenominator & \AFDFalseBlockingRate \\
Expected-action mismatch outside the primary denominators & \AFDExpectedActionMismatches{} / pending & pending \\
Reason mismatch & \AFDReasonMismatches{} / pending & pending \\
Pipeline error & \AFDPipelineErrors{} / \AFDCorpusSamples & pending \\
\bottomrule
\end{tabularx}
\end{table}

Generated tables repeat counts by stratum, concrete format, carrier mechanism,
policy oracle, and reconstruction requirement where the denominator is nonzero.
Ground-truth construction fields remain distinct from policy-oracle fields.
Delivery of an expected block is false acceptance. Withholding an expected-clean
S0 or S1 record is primary benign false blocking. Retention of a named carrier
property after an expected-clean reconstruction is a separate endpoint.

No row is described as malware-family detection accuracy. S6 concerns
unsupported carrier classes. S3 concerns evidence disagreement. S4 concerns
canonical representation and complete consumption. S7 contains named
regressions only. A behavioral label does not change the measured system action.

\subsection{Invariant, differential, and resilience evidence}
\label{subsec:results-properties}

The T1--T8 report provides eligible cases, passes, failures, exclusions, and
the corresponding F1--F8 cross-reference. Canonical failures are
\AFDCanonicalGateFailures. A missing required test is not interpreted as a
pass. It makes the relevant release gate false.

Second-pass results are reported outside T1--T8. They split
\AFDExactStableOutputs exact deterministic outputs from
\AFDSemanticClosureOutputs lossy semantic-closure outputs. Lossy profiles also
report the separately frozen fidelity-drift comparator. These categories are
not combined under the label idempotence.

Differential validation reports accepted outputs, evaluated outputs, parser
families, disagreements, tool errors, and \AFDDifferentialFailures failures.
Fidelity reports metric profiles, denominators, thresholds, tool errors, and
\AFDFidelityFailures failures. Fuzzing reports every repetition, seed, budget,
execution count, coverage observation, and \AFDFuzzFailures unresolved failure.
Fault injection reports every planned and implemented injection, including
\AFDFaultFailures violations of \cref{eq:no-publication}. An absent report never
creates a zero placeholder.

\subsection{Performance and fidelity presentation}
\label{subsec:results-performance}

End-to-end warm median latency is \AFDMedianLatency and p95 latency is
\AFDPninetyfiveLatency. Their hierarchical interval field is
\AFDPerformanceCI. Sequential single-worker service rate is \AFDServiceRate,
and peak resident memory is \AFDPeakMemory. The generated table stratifies
these values by profile, format, reconstruction level, size band, and accepted
or rejected path. Cold and warm modes remain separate. Counts for the ten
sessions, warm-ups, observations, timeouts, base artifacts, and 10,000 bootstrap
replicates accompany each summary.

Fidelity is summarized as \AFDFidelitySummary only after each metric profile,
denominator, threshold, and failure is present. Security and fidelity tables
remain separate. Canonical acceptance cannot conceal unusable media, and high
media quality cannot compensate for a security-gate failure.

\subsection{Append-only regeneration and interpretation}
\label{subsec:results-regeneration}

Raw JSON Lines, adapter reports, and hashed scripts regenerate every table and
macro for \AFDRunID. The run registry is append-only. A rerun never overwrites
its predecessor and must cite the predecessor identifier and reason. Every
interpretation names the run, stratum, denominator, profile, reconstruction
level, and gate value. Until a generated run artifact exists, all outcomes
remain pending.
